\begin{problem}[2020第37届全国中学生物理竞赛复赛][双黑洞旋近引力波辐射]{123}
    激光干涉引力波天文台（LIGO）2015年首次探测到了十亿光年外双黑洞并合产生的引力波，证实了爱因斯坦理论关于存在引力波和黑洞的预言。黑洞并合过程分为三个阶段：第一阶段（旋近阶段），两黑洞围绕系统质心在同一平面内做近似圆周的运动，损失的机械能转化为引力辐射能，两者螺旋式逐渐靠近；第二阶段（并合阶段），双黑洞并合为一个黑洞；最后阶段（衰荡阶段），并合的黑洞弛豫至平衡状态，成为一个稳定的旋转黑洞。在旋近阶段，若忽略黑洞的自转和大小，则双黑洞均可视为质量分别恒为 $M$ 和 $m$ 的质点，它们之间距离 $L$ 随时间而逐渐减小。假定系统除了辐射引力波外无其它能量耗散，不考虑引力辐射的反作用，可用牛顿引力理论进行近似处理。
    \begin{subproblem}
        引力波辐射功率除了与引力常量 $G$ 成正比之外，还可能与两黑洞的质量 $M$ 与 $m$、两黑洞之间的距离 $L$、以及系统绕质心的转动惯量 $I$、转动角速度 $\omega$ 和辐射引力波的传播速度（其大小等于真空中的光速 $c$）有关，试选取描述转动体系辐射的三个物理量与 $G$ 一起导出引力波辐射功率的表达式（假设其中可能待定的无量纲比例常数为 $\alpha$）；
    \end{subproblem}
    \begin{subproblem}
        若在初始时刻 $t = 0$ 时两黑洞之间的距离为 $L_{0}$，且引力波辐射功率表达式中的无量纲比例常数 $\alpha$ 是已知的，求两黑洞从 $t = 0$ 时开始绕系统质心旋转 $360^{\circ}$ 所需要的时间；
    \end{subproblem}
    \begin{subproblem}
        当两黑洞从 $t = 0$ 时开始绕系统质心旋转多少度时，它们间的距离恰好是其初始距离的一半？
    \end{subproblem}
\end{problem}